\documentclass[12pt]{article}
\usepackage[utf8]{inputenc}
\usepackage[T1]{fontenc}
\usepackage[margin=1in]{geometry}
\usepackage{setspace}
\usepackage{xcolor}
\usepackage[colorlinks=true, urlcolor=blue!60!black, linkcolor=blue!60!black]{hyperref}
\setlength{\parindent}{2em}
\setlength{\parskip}{1.25em}
\renewcommand{\baselinestretch}{1.0}
\begin{document}
\begin{center}
{\LARGE \textit{Lochner v. New York}, 198 U.S. 45 (1905)\par}
\vspace{0.8em}
{\large Case Brief --- Four-Page Edition\par}
\end{center}
\vspace{1em}

\section{Memory Jogger}

A New York law capping bakers' hours at ten per day and sixty per week violated the ``liberty of contract'' protected by the Due Process Clause, because it lacked a direct and reasonable relation to public health or worker safety (\href{https://caselaw.findlaw.com/court/us-supreme-court/198/45.html}{FindLaw}).

\section{Detailed Case Facts}

New York's Bakeshop Act regulated bakeries' drainage, plumbing, ventilation, ceilings, washrooms, and inspection (Sections~111--115). Section~110 separately barred any employee from being required or permitted to work in a bakery more than sixty hours a week or ten hours a day, with no emergency exception. Joseph Lochner, who ran a small bakery in Utica and had already been fined \$25 for a first violation, was indicted in 1901 for permitting an employee, Aman Schmitter, to work more than sixty hours in a week. As the courts treated the record, the employee agreed to the extra hours to earn more wages.

\section{Procedural History}

The Oneida County Court overruled Lochner's demurrer, convicted him of a second-offense misdemeanor, and fined him \$50. The Appellate Division affirmed 3--2 and the New York Court of Appeals affirmed 4--3, sustaining Section~110 as a health law. The Supreme Court heard the case on writ of error.

\section{Judicial Votes}

Reversed, \textbf{5--4} (decided April 17, 1905).
\begin{itemize}
\item \textbf{Majority:} Justice Peckham, joined by Chief Justice Fuller and Justices Brewer, Brown, and McKenna.
\item \textbf{Dissents:} Justice Harlan, joined by Justices White and Day; Justice Holmes, alone.
\end{itemize}

\section{Holding}

The Due Process Clause protects a general liberty of contract, including the right to buy and sell labor. A state may restrict it through the police power, but a law interfering with labor contracts between competent adults is invalid unless it bears a direct relation, as a means to an end, to a legitimate police-power objective and rests on a fair, reasonable showing of material danger to health. Section~110 is not a valid health law. Peckham's approval of the sanitary sections and his warnings about capping the hours of clerks and doctors are dicta.

\section{Analysis of Opinions}

\textbf{Majority (Peckham).} Relying on \textit{Allgeyer v. Louisiana}, Peckham treated Section~110 as an absolute bar on willing parties contracting for overtime. He distinguished \textit{Holden v. Hardy} (miners, an acutely hazardous trade, with an emergency exception), \textit{Atkin v. Kansas} (public works), and \textit{Jacobson v. Massachusetts} (vaccination). Courts must decide whether a law is a fair exercise of the police power or an arbitrary interference, lest the police power become a ``mere pretext.'' As a labor law, the act failed because bakers are competent adults; as a public-health law, the wholesomeness of bread does not depend on bakers' hours; and as a worker-health law, baking was not unusually unhealthy, and accepting the rationale would leave every occupation ``at the mercy of legislative majorities.'' The separate sanitary sections showed the hours cap was passed ``from other motives.''

\textbf{Harlan (dissenting).} Accepting liberty of contract, Harlan insisted on a strong presumption of constitutionality. Medical treatises and labor statistics showed that bakers suffered from flour dust, overheated night work, and shortened lives; because there was ``room for debate and for an honest difference of opinion,'' the Court should defer.

\textbf{Holmes (dissenting).} The case was decided ``upon an economic theory which a large part of the country does not entertain.'' Sunday laws, usury laws, and vaccination mandates already restrict contract. Liberty is perverted when used to block the natural outcome of a dominant opinion unless a rational and fair man would necessarily see a violation of fundamental principles.

\section{Examples: Future Applications}

\begin{enumerate}
\item \textbf{Same-Side (Invalid): Retail Clerks.} A 48-hour weekly cap for adult store clerks rests on general fatigue, not a direct health nexus.
\item \textbf{Same-Side (Invalid): Minimum Wage.} A minimum wage for adult restaurant workers is pure labor regulation (as in \textit{Adkins}, later overruled by \textit{West Coast Hotel}).
\item \textbf{Opposite-Side (Valid): Underground Mines.} An eight-hour day for sulfur-mine workers with an emergency exception follows \textit{Holden v. Hardy}.
\item \textbf{Opposite-Side (Valid): Bakery Sanitation.} Exhaust fans, tile floors, and fire exits directly protect health, like Sections~111--115.
\item \textbf{Fence-Sitter: Bus Drivers.} A ten-hour cap for intercity bus drivers regulates a common occupation, but a fatigued driver directly endangers passengers and the public---unlike a tired baker's bread.
\end{enumerate}

\section{Critique}

\textbf{Progressive critiques.} Critics from the Progressive Era onward condemned \textit{Lochner} as judges substituting laissez-faire preferences for democratic choices. Sunstein located its error in treating common-law baselines as natural and neutral; Gillman traced it to an anti-class-legislation tradition that proved unworkable in an industrial economy. \textit{Nebbia} (1934) and \textit{West Coast Hotel} (1937) repudiated it.

\textbf{Originalist and libertarian views.} Bork and Scalia condemn substantive due process itself, tracing \textit{Roe} to \textit{Lochner}. Revisionists such as Bernstein reject the Social Darwinism narrative, noting that the Bakeshop Act favored unionized large bakeries over immigrant-owned shops, and Barnett and Bernick argue that due process properly bars arbitrary legislation.

\section{Key Quotations}

\begin{quote}
``General propositions do not decide concrete cases. The decision will depend on a judgment or intuition more subtle than any articulate major premise.''
\end{quote}
\noindent --- Justice Holmes, dissenting, 198 U.S. at 76.

\begin{quote}
``The Fourteenth Amendment does not enact Mr. Herbert Spencer's \textit{Social Statics}. ... But a constitution is not intended to embody a particular economic theory, whether of paternalism and the organic relation of the citizen to the state or of \textit{laissez faire}. It is made for people of fundamentally differing views, and the accident of our finding certain opinions natural and familiar or novel and even shocking ought not to conclude our judgment upon the question whether statutes embodying them conflict with the Constitution of the United States.''
\end{quote}
\noindent --- Justice Holmes, dissenting, 198 U.S. at 75--76.

\begin{quote}
``The mere assertion that the subject relates, though but in a remote degree, to the public health, does not necessarily render the enactment valid. The act must have a more direct relation, as a means to an end, and the end itself must be appropriate and legitimate, before an act can be held to be valid which interferes with the general right of an individual to be free in his person and in his power to contract in relation to his own labor.''
\end{quote}
\noindent --- Justice Peckham, for the Court, 198 U.S. at 57--58.

\end{document}
